% ===================================================================
%  Dodatek I: Stabilizacja potencjału i trzecia generacja leptonów
%  Mechanizm V_mod = V_orig + lambda*(psi-1)^6/6
%  Wyniki numeryczne z poprawioną kwadraturą (siatka log)
%  TGP v1 — 2026-03-23
% ===================================================================
\section{Stabilizacja potencjału TGP i trzecia generacja cząstek}%
\label{app:v2-potencjal}

\subsection{Motywacja: problem z katastrofą kwartyką}%
\label{app:I-motywacja}

Potencjał TGP w wersji wyjściowej~\eqref{eq:Veff-kink}:
\begin{equation}\label{eq:V-orig}
    V_{\mathrm{orig}}(\psi)
    = \frac{\psi^3}{3} - \frac{\psi^4}{4},
    \qquad \psi = \frac{\Phi}{\PhiZero},
\end{equation}
posiada lokalny ekstrema w~$\psi = 0$ i~$\psi = 1$ (próżnia TGP).
Dla $\psi > 1$ potencjał maleje monotonicznie:
$V_{\mathrm{orig}}(\psi) \sim -\psi^4/4 \to -\infty$.

Warunek samospójności cząstki jako solitonu Yukawa definiuje
\textbf{funkcję g}:
\begin{equation}\label{eq:g-def}
    g(K) \;\equiv\; \frac{E(K)}{K} - 4\pi \;\stackrel{!}{=}\; 0,
\end{equation}
gdzie $K$ jest amplitudą profilu Yukawa~$\Phi(r) = \PhiZero + K e^{-m_{\rm sp}r}/r$,
a $E(K)$ — całkowita energia solitonu. Rozwiązania $g(K_i) = 0$
odpowiadają samospójnym masom cząstek.

\begin{proposition}[Dwie generacje w potencjale oryginalnym]%
\label{prop:dwa-zera}
Dla~$V_{\mathrm{orig}}$ funkcja $g(K)$ posiada \emph{dokładnie dwa zera}:
$K_1$ (mała amplituda, reżim liniowy) i~$K_2$ (duża amplituda, reżim
kwartyczny). Stosunek~$r_{21} = K_2/K_1 \approx 207$
odpowiada stosunkowi mas mionu do elektronu.
Dla $K > K_2$: $g(K) \to -\infty$ monotonicznie —
brak trzeciego zera, brak trzeciej generacji.
\end{proposition}

\begin{proof}[Uzasadnienie analityczne]
Dla dużych $K$ rdzeń solitonu osiąga $\psi_{\rm core} \gg 1$.
Wkład potencjalny:
\[
    E_{\rm pot} \approx 4\pi \int \frac{-\psi_{\rm core}^4}{4}\, r^2\, dr
    \;\sim\; -\frac{K^4}{\PhiZero^2\, a_\Gamma} \;\to\; -\infty.
\]
Brak członu stabilizującego rosnącego szybciej niż $\psi^4$
sprawia, że $g(K) \to -\infty$ i~brak dalszych zer.
\end{proof}

% -------------------------------------------------------------------
\subsection{Modyfikacja potencjału: człon stabilizujący}%
\label{app:I-Vmod}
% -------------------------------------------------------------------

Proponujemy minimalne rozszerzenie potencjału o~człon szóstego stopnia:
\begin{equation}\label{eq:V-mod}
    \boxed{V_{\rm mod}(\psi)
    = \frac{\psi^3}{3} - \frac{\psi^4}{4}
    + \frac{\lambda\,(\psi-1)^6}{6}},
    \qquad \lambda > 0.
\end{equation}

\begin{proposition}[Wyprowadzenie członu $(\psi-1)^6$ z~renormalizacji substratu]%
\label{prop:V6-from-substrate}
\statuslabel{Propozycja $\to$ Twierdzenie (1-loop)}
Cz\l{}on $\lambda(\psi-1)^6/6$ nie jest \emph{ad hoc}, lecz wynika
z~jednoparametrowej renormalizacji Wilsona hamiltonianu
substratu~$H_\Gamma$.
Poni\.zszy argument jest pe\l{}ny na poziomie 1-loop;
pe\l{}ne twierdzenie wymaga nieperturbacyjnych danych MC
(weryfikacja: \texttt{scripts/wilson\_rg\_vmod.py}).

\medskip\noindent\textbf{Argument.}
Efektywny potencjał Ginzburga--Landaua substratu $\Gamma$
(3D siatka z~symetri\k{a}~$\mathbb{Z}_2$) ma postać:
\begin{equation}\label{eq:GL-substrate}
    F[v] = \int\!\Bigl[\tfrac{1}{2}(\nabla v)^2
    + \tfrac{r}{2}\,v^2 + \tfrac{u_4}{4!}\,v^4
    + \tfrac{u_6}{6!}\,v^6 + \cdots\Bigr]d^3x,
\end{equation}
gdzie $v(x)$ jest polem porz\k{a}dku ($v^2 \propto \langle\hat{s}_i^2\rangle
\propto \Phi$).
Symetria $\mathbb{Z}_2$ ($v \to -v$) wyklucza nieparzyste potęgi.

\medskip\noindent\emph{Krok~1: Renormalizacja Wilsona.}\;
W~punkcie stałym Wilsona--Fishera (WF) w~$d = 3$ wymiarach
funkcja beta sprzężenia $u_6$ ma postać%
\footnote{Zinn-Justin, \emph{Quantum Field Theory and Critical Phenomena},
Chap.~25; Pelissetto \& Vicari, \emph{Phys.\ Rep.}\ 368 (2002) 549.}:
\begin{equation}\label{eq:beta-u6}
    \beta_{u_6} = \epsilon\, u_6 + c_1\, u_4^2 + c_2\, u_4\, u_6 + \cdots,
    \qquad \epsilon = 4 - d = 1.
\end{equation}
Nawet je\'sli $u_6^{\rm bare} = 0$ na skali UV (substrat dyskretny),
człon $c_1 u_4^2$ \textbf{generuje} $u_6^* \neq 0$ w~punkcie stałym WF.
Konkretnie:
\begin{equation}\label{eq:u6-WF}
    u_6^* = -\frac{c_1}{c_2}\, u_4^* + O(u_4^{*2}).
\end{equation}

\medskip\noindent\emph{Krok~2: Mapowanie na zmienn\k{a} $\psi$.}\;
Pole porz\k{a}dku substratu $v$ jest zwi\k{a}zane z~polem TGP przez
$\Phi = v^2 / \Phi_0$, tj.~$\psi = v^2 / v_0^2$,
gdzie $v_0 = \sqrt{\Phi_0}$ jest warto\'sci\k{a} pr\'o\.zniow\k{a}.
Fluktuacja $\delta v = v - v_0$ wyra\.za si\k{e} przez $\psi$
\textbf{nieliniowo}:
\begin{equation}\label{eq:dv-to-psi}
    \frac{\delta v}{v_0} = \sqrt{\psi} - 1
    = \frac{\psi - 1}{2} - \frac{(\psi-1)^2}{8} + \cdots
\end{equation}
Cz\l{}on $u_6^*\,v^6 / 6!$ w~funkcjonale GL~\eqref{eq:GL-substrate}
rozwija si\k{e} wok\'o\l{} pr\'o\.zni:
\begin{equation}\label{eq:F-to-Vmod}
    \frac{u_6^*}{6!}\,(\delta v)^6
    = \frac{u_6^*\,v_0^6}{6!}\,(\sqrt{\psi}-1)^6
    = \frac{u_6^*\,v_0^6}{6!\cdot 2^6}\,(\psi-1)^6
      \bigl[1 + O(\psi-1)\bigr].
\end{equation}
Identyfikuj\k{a}c z~cz\l{}onem $\lambda(\psi-1)^6/6$
w~\eqref{eq:V-mod}:
\begin{equation}\label{eq:lambda-from-WF}
    \lambda_{\rm eff}
    = \frac{|u_6^*|\,v_0^6}{5!\cdot 2^6}
    = \frac{|u_6^*|\,\PhiZero^3}{5!\cdot 64}.
\end{equation}
\textbf{Oszacowanie numeryczne.}\;
Sprzężenie $u_6^*$ jest bezwymiarowe na skali odci\k{e}cia
$\Lambda = 1/a_\Gamma$.
W~konwencji TGP (potencja\l{} $V(\psi)$ bezwymiarowy)
efektywne $\lambda$ jest t\l{}umione przep\l{}ywem RG
od skali substratowej do skali solitonu.
Z~analizy wymiarowej:
\begin{equation}\label{eq:lambda-dimensional}
    \lambda_{\rm eff}
    \sim \frac{a_\Gamma^2}{\PhiZero^2}
    \approx \frac{(0{,}040)^2}{(24{,}66)^2}
    \approx 2{,}6 \times 10^{-6},
\end{equation}
co jest sp\'ojne z~numerycznym przedzia\l{}em
$\lambda \sim 10^{-7}\text{--}10^{-5}$ wymaganym przez trzy generacje
(weryfikacja: \texttt{scripts/wilson\_rg\_vmod.py}).
Przy hipotezie $a_\Gamma\cdot\PhiZero \approx 1$ mamy
$\lambda \sim 1/\PhiZero^4 \approx 2{,}7 \times 10^{-6}$
--- identyczny rz\k{a}d wielko\'sci.

\medskip\noindent\emph{Krok~3: Dlaczego akurat stopie\'n 6, nie 8?}\;
Cz\l{}on $u_8^* v^8 / 8!$ jest r\'ownie\.z generowany przez RG,
ale jest \textbf{podw\'ojnie t\l{}umiony}:
\begin{enumerate}[nosep]
  \item \emph{Perturbacyjnie}: $u_8^* \sim u_4^{*3}$ (diagram 3-loop),
    podczas gdy $u_6^* \sim u_4^{*2}$ --- t\l{}umienie o~czynnik $u_4^*$.
  \item \emph{Wymiar anomalny}: wymiar kanoniczny $u_8$ wynosi
    $\Delta_8 = 2\epsilon$ (vs.\ $\Delta_6 = \epsilon$), wi\k{e}c
    efektywne $\lambda_8$ jest t\l{}umione dodatkowym czynnikiem
    $a_\Gamma^2$ wzgl\k{e}dem $\lambda_6$:
    $\lambda_8/\lambda_6 \sim a_\Gamma^2/\PhiZero^2 \approx 10^{-6}$.
\end{enumerate}
Na skali generacji (solitony o~$\psi_{\rm core} \lesssim 10^3$)
cz\l{}on $\psi^8$ jest zaniedbywalny, co uzasadnia obci\k{e}cie
do stopnia 6.
\end{proposition}

\begin{proposition}[Warunki próżni zachowane]%
\label{prop:proznia-zachowana}
Modyfikacja~\eqref{eq:V-mod} zachowuje strukturę próżni TGP:
\begin{enumerate}[nosep,label=(\roman*)]
    \item $V_{\rm mod}(1) = V_{\rm orig}(1) = 1/12$
          \quad (energia próżni niezmieniona),
    \item $V'_{\rm mod}(1) = 0$
          \quad (warunek równowagi — próżnia w~$\psi=1$),
    \item $V''_{\rm mod}(1) = -1$
          \quad (masa skalaru niezmieniona),
    \item $V_{\rm mod}(\psi) \to +\infty$ dla $\psi \to +\infty$
          \quad (potencjał ograniczony od dołu).
\end{enumerate}
\end{proposition}

\begin{proof}
Pochodne członu stabilizującego $\delta V = \lambda(\psi-1)^6/6$:
$\delta V(1) = 0$, $\delta V'(1) = 0$, $\delta V''(1) = 0$
(sześć pochodnych znika w~$\psi=1$). Własności (i)--(iii) wynikają
wprost. Własność (iv): $\lambda\psi^6/6 > \psi^4/4$ gdy
$\psi^2 > 3/(2\lambda) \equiv \psi_{\rm cross}^2$;
dla $\psi > \psi_{\rm cross}$ człon $\lambda$-owy dominuje.
\end{proof}

\begin{remark}[Nowe minimum i~głębokość studni]%
\label{rem:nowe-minimum}
Potencjał $V_{\rm mod}$ posiada globalne minimum przy
\begin{equation}\label{eq:psi-new}
    \psi_{\rm new} \approx \frac{1}{\sqrt{\lambda}},
\end{equation}
z~głębokością $V_{\rm mod}(\psi_{\rm new}) - V_{\rm mod}(1) \approx -\lambda^{-2}/4$
(rząd $10^{11}$ dla typowych $\lambda \sim 10^{-6}$).
Głębokość ta \emph{nie ma znaczenia fizycznego} dla solitonów
z~rdzeniem $\psi_{\rm core} \ll \psi_{\rm new}$ (generacje 1 i~2),
lecz wpływa na geometrię funkcji $g(K)$ w~pobliżu trzeciego zera
(patrz twierdzenie~\ref{thm:trzy-zera}).
\end{remark}

% -------------------------------------------------------------------
\subsection{Topologia funkcji~$g(K)$: skąd biorą się trzy generacje}%
\label{app:I-topologia}
% -------------------------------------------------------------------

Kluczowy wynik analityczny: liczba zer $g(K)$ jest bezpośrednio
związana z~liczbą \emph{przekroczeń} poziomej linii $V_{\rm mod}(\psi) = V_{\rm mod}(1)$
przez potencjał dla $\psi > 1$.

\begin{theorem}[Trzy zera $g(K)$ dla potencjału stopnia 6]%
\label{thm:trzy-zera}
Niech $\psi_{\rm cross}$ będzie jedynym rozwiązaniem
$V_{\rm mod}(\psi) = V_{\rm mod}(1)$ dla $\psi > 1$:
\begin{equation}\label{eq:psi-cross}
    \psi_{\rm cross} \approx \sqrt{\frac{3}{2\lambda}}
    = \sqrt{\frac{3}{2}}\cdot\psi_{\rm new}.
\end{equation}
Wówczas:
\begin{itemize}[nosep]
    \item Dla $1 < \psi < \psi_{\rm cross}$:
          $V_{\rm mod}(\psi) < V_{\rm mod}(1)$
          (wkład potencjalny rdzenia \emph{ujemny}),
    \item Dla $\psi > \psi_{\rm cross}$:
          $V_{\rm mod}(\psi) > V_{\rm mod}(1)$
          (wkład potencjalny rdzenia \emph{dodatni}).
\end{itemize}
Istnienie jednego przekroczenia $\psi_{\rm cross}$ tworzy jeden
dodatkowy garb w~funkcji $g(K)$ (poza dwiema gałęziami istniejącymi
dla $\lambda = 0$), dając łącznie \emph{dokładnie trzy zera}.
Dla $K > K_3$: rdzeń solitonu przekracza $\psi_{\rm cross}$,
powłoka zewnętrzna (\emph{shell}) dominuje energetycznie
i~$g(K) \to -\infty$ — brak czwartego zera.
\end{theorem}

\begin{corollary}[Liczba generacji a~stopień potencjału]%
\label{cor:stopien}
Potencjał stopnia $2n$ (z~członem stabilizującym $\sim\psi^{2n}$,
$n \geq 3$) generuje co najwyżej $(n-1)$ zer funkcji $g(K)$,
a~zatem co najwyżej $(n-1)$ samospójnych generacji cząstek.
Dla stopnia 6 ($n=3$): \emph{maksymalnie trzy generacje}.

Fakt eksperymentalny — istnienie dokładnie trzech generacji leptonów —
jest spójny z~potencjałem stopnia 6 i~\emph{predykuje}, że potencjał
TGP nie zawiera członów $\lambda'\psi^8$, $\lambda''\psi^{10}$, itd.
\end{corollary}

% -------------------------------------------------------------------
\subsection{Problem techniczny: kwadryatura i~siatka logarytmiczna}%
\label{app:I-kwadryatura}
% -------------------------------------------------------------------

\begin{remark}[Krytyczny błąd wcześniejszej analizy]%
\label{rem:blad-kwadratury}
Wcześniejsze obliczenia (skrypt \texttt{v2\_fine\_scan.py},
równomierna siatka $N = 800$) dawały $K_3 \approx 47.3$,
$r_{31} \approx 3477$ dla $\lambda^* \approx 5.51 \times 10^{-7}$.
Były to \emph{błędne wyniki} — artefakt niedostatecznej kwadratury,
\textbf{nie} rozbieżność fizyczna.
\end{remark}

Całka energetyczna $E(K) = \int [\text{Ek} + \text{Ep}]\, r^2\, dr$
dla dużych $K$ (rdzeń $\psi_{\rm core} \gg 1$) zawiera
\emph{ostre piecie} integrandu przy $r = a_\Gamma$:
czynnik $\lambda(\psi-1)^6 \cdot r^2$ maleje od wartości
$\sim \lambda \psi_{\rm core}^6 a_\Gamma^2$ do zera na skali
$r_{\rm cross} = K / \psi_{\rm cross} \gtrsim a_\Gamma$.

Równomierna siatka o~$N$ punktach od $a_\Gamma$ do $r_{\rm max}$
daje znikomą rozdzielczość przy $r \approx a_\Gamma$
i~wymaga $N \to \infty$ do zbieżności.
Rozwiązanie: \textbf{siatka logarytmiczna}:
\begin{equation}\label{eq:log-grid}
    r_k = a_\Gamma \left(\frac{r_{\rm max}}{a_\Gamma}\right)^{k/N},
    \qquad k = 0, 1, \ldots, N,
\end{equation}
która zagęszcza punkty przy $r \approx a_\Gamma$, zapewniając
zbieżność już przy $N \sim 500$.

\begin{table}[h]
\centering
\small
\caption{Porównanie zbieżności $E(K)/K$ dla $K = 47.3$,
  $\alpha = 7.0$, $a_\Gamma = 0.030$, $\lambda = 5.514\times10^{-7}$
  (skrypt \texttt{v2\_kwadryatura.py}).}
\label{tab:kwadryatura-porownanie}
\begin{tabular}{@{}lccc@{}}
\toprule
\textbf{Siatka} & $N=500$ & $N=8000$ & \textbf{Wniosek} \\
\midrule
Równomierna    & $+1.67\times10^{9}$ & $-1.05\times10^{8}$ & Zmiana znaku! \\
\textbf{Logarytmiczna} & $\mathbf{-2.39\times10^{8}}$
                       & $\mathbf{-2.39\times10^{8}}$
                       & \textbf{Zbieżna natychmiast} \\
Adaptywna (\texttt{quad}) & \multicolumn{2}{c}{$-2.39\times10^{8}$}
                          & Potwierdza log \\
\bottomrule
\end{tabular}
\end{table}

Potwierdzenie pochodzi z~adaptywnej kwadratury
\texttt{scipy.integrate.quad} (bez żadnej siatki dyskretnej),
która daje identyczny wynik co siatka logarytmiczna.
\textbf{Poprzednia ``rozbieżność'' była w~100\% artefaktem
siatki równomiernej.}

% -------------------------------------------------------------------
\subsection{Dowód: trzeci soliton jest rzeczywisty}%
\label{app:I-dowod}
% -------------------------------------------------------------------

\begin{theorem}[Rzeczywistość trzeciego zera]%
\label{thm:trzecie-zero}
Dla parametrów $\alpha = 7.0$, $a_\Gamma = 0.030$,
$\lambda = 5.514 \times 10^{-7}$ funkcja $g(K)$ posiada
\emph{prawdziwe} trzecie zero w~przedziale $K \in (80, 82)$,
tzn.\ zmiana znaku jest zbieżna i~niezależna od $N$.
\end{theorem}

\begin{proof}[Weryfikacja numeryczna (skrypt \texttt{v2\_konwergencja\_finalna.py})]
Tabela~\ref{tab:konwergencja-g} zestawia wartości $g(K)$
dla $K = 80$ i~$K = 82$ przy rosnącym $N$ (siatka logarytmiczna):

\begin{table}[h]
\centering
\small
\caption{Zbieżność $g(K) = E(K)/K - 4\pi$ z~$N$
  (siatka log, $\lambda = 5.514\times10^{-7}$, $\alpha=7.0$, $a_\Gamma=0.030$).}
\label{tab:konwergencja-g}
\begin{tabular}{@{}rrrr@{}}
\toprule
$N$ & $g(K=80)$ & $g(K=82)$ & $\Delta$ [od poprz.] \\
\midrule
1\,000  & $-9.9869\times10^{5}$ & $+8.8720\times10^{5}$ & — \\
2\,000  & $-9.9862\times10^{5}$ & $+8.8718\times10^{5}$ & $0.007\%$ \\
4\,000  & $-9.9861\times10^{5}$ & $+8.8717\times10^{5}$ & $0.001\%$ \\
8\,000  & $-9.9860\times10^{5}$ & $+8.8717\times10^{5}$ & $<0.001\%$ \\
16\,000 & $-9.9860\times10^{5}$ & $+8.8717\times10^{5}$ & $<0.001\%$ \\
Adaptywna & $-9.9860\times10^{5}$ & $+8.8717\times10^{5}$ & Potwierdza \\
\bottomrule
\end{tabular}
\end{table}

Zmiana znaku $g(80) < 0 < g(82)$ jest stabilna:
relatywna zmiana z~$N=1000$ do $N=16000$ wynosi $<0.01\%$.
Adaptywna kwadryatura potwierdza wynik siatki logarytmicznej.
Twierdzenie Brentq-a gwarantuje istnienie zera $K_3 \in (80,82)$;
bisekcja daje $K_3 \approx 81.0$.

\emph{Poprzednia analiza stwierdzająca rozbieżność całki
była błędem kwadratury, nie cechą fizyczną rozwiązania.}
\end{proof}

% -------------------------------------------------------------------
\subsection{Wyznaczenie parametrów: skan 2D}%
\label{app:I-skan2d}
% -------------------------------------------------------------------

Cel: znaleźć trójkę $(\alpha, a_\Gamma, \lambda)$ taką, że
\begin{equation}\label{eq:cel}
    r_{21} = \frac{K_2}{K_1} = 207, \qquad r_{31} = \frac{K_3}{K_1} = 3477.
\end{equation}

\textbf{Metoda} (skrypty \texttt{v2\_2d\_skan\_log.py},
\texttt{v2\_bisekcja\_alpha.py}):
\begin{enumerate}[nosep]
    \item Dla każdego punktu $(\alpha, a_\Gamma)$ w~siatce:
          bisekcja po $\lambda$ wyznacza $\lambda^*(\alpha, a_\Gamma)$
          takie, że $r_{31} = 3477$ (metoda Brentq, tolerancja $10^{-3}$).
    \item Przy wyznaczonym $\lambda^*$ oblicza się $r_{21}(\alpha, a_\Gamma)$.
    \item Dla ustalonego $a_\Gamma$: bisekcja po $\alpha^*$ wyznacza
          wartość dającą $r_{21} = 207$.
\end{enumerate}

Układ dwóch równań~\eqref{eq:cel} w~trzech nieznanych
$(\alpha, a_\Gamma, \lambda)$ ma rozwiązania tworzące
\emph{jednoparametrową krzywą} w~przestrzeni parametrów.

\begin{table}[h]
\centering
\small
\caption{Jednoparametrowa rodzina rozwiązań równań~\eqref{eq:cel}
  (skrypt \texttt{v2\_bisekcja\_alpha.py}, siatka log $N=1500$).}
\label{tab:rodzina-rozwiazan}
\begin{tabular}{@{}ccccccccc@{}}
\toprule
$\alpha^*$ & $a_\Gamma$ & $\lambda^*$ & $K_1$ & $K_2$ & $K_3$
    & $r_{21}$ & błąd $r_{21}$ & $r_{31}$ \\
\midrule
$5.9148$ & $0.025$ & $2.883\times10^{-6}$
    & $0.00853$ & $1.76610$ & $29.670$ & $206.972$ & $-0.013\%$ & $3477.1$ \\
$6.8675$ & $0.030$ & $3.743\times10^{-6}$
    & $0.00896$ & $1.85472$ & $31.160$ & $206.969$ & $-0.015\%$ & $3477.1$ \\
$7.7449$ & $0.035$ & $4.618\times10^{-6}$
    & $0.00939$ & $1.94375$ & $32.655$ & $206.967$ & $-0.016\%$ & $3477.1$ \\
$\mathbf{8.5616}$ & $\mathbf{0.040}$ & $\mathbf{5.501\times10^{-6}}$
    & $\mathbf{0.00982}$ & $\mathbf{2.03272}$ & $\mathbf{34.144}$
    & $\mathbf{206.999}$ & $\mathbf{-0.000\%}$ & $\mathbf{3477.0}$ \\
\bottomrule
\end{tabular}
\end{table}

\begin{remark}[Najlepsze rozwiązanie]
Punkt $(\alpha^*, a_\Gamma) = (8.5616, 0.040)$ z~$\lambda^* = 5.501\times10^{-6}$
daje $r_{21} = 206.999$ (błąd $-0.000\%$) i~$r_{31} = 3477.0$ (błąd $0.0\%$).
Jest to rozwiązanie \emph{praktycznie dokładne} przy 4-cyfrowej precyzji.
\end{remark}

\textbf{Weryfikacja stabilności} (skrypt \texttt{v2\_weryfikacja\_finalna.py}):
\begin{table}[h]
\centering
\small
\caption{Stabilność zer $K_i$ względem $N$ dla $\alpha=8.5616$,
  $a_\Gamma=0.040$, $\lambda^*=5.501\times10^{-6}$.}
\label{tab:stabilnosc-zer}
\begin{tabular}{@{}rccccc@{}}
\toprule
$N$ & $K_1$ & $K_2$ & $K_3$ & $r_{21}$ & $r_{31}$ \\
\midrule
$2000$ & $0.009820$ & $2.032726$ & $34.14423$ & $206.999$ & $3477.0$ \\
$4000$ & $0.009820$ & $2.032728$ & $34.14444$ & $206.998$ & $3477.0$ \\
$8000$ & $0.009820$ & $2.032728$ & $34.14450$ & $206.998$ & $3477.0$ \\
\bottomrule
\end{tabular}
\end{table}

Całkowita zmiana $K_3$ od $N=2000$ do $N=8000$: $\Delta K_3 = 0.00027$,
co stanowi $0.0008\%$ wartości. Wszystkie trzy zera są \emph{numerycznie stabilne}.

% -------------------------------------------------------------------
\subsection{Universalna relacja rdzenia i~minimum potencjału}%
\label{app:I-relacja}
% -------------------------------------------------------------------

\begin{proposition}[Universalna relacja $\psi_{\rm core}/\psi_{\rm new}$]%
\label{prop:universalna-relacja}
Dla każdego punktu z~rodziny rozwiązań (tabela~\ref{tab:rodzina-rozwiazan})
stosunek wartości pola w~rdzeniu trzeciego solitonu do położenia
nowego minimum potencjału spełnia:
\begin{equation}\label{eq:universalna-relacja}
    \frac{\psi_{\rm core}(M_3)}{\psi_{\rm new}}
    = \frac{1 + K_3\, e^{-a_\Gamma}/a_\Gamma}{1/\sqrt{\lambda^*}}
    \;\in\; [1.926,\; 1.967]
    \;\approx\; 2,
\end{equation}
mimo że poszczególne parametry $(\alpha, a_\Gamma, \lambda^*)$
zmieniają się ponad dwukrotnie wzdłuż krzywej rozwiązań.
\end{proposition}

\begin{table}[h]
\centering
\small
\caption{Stosunek $\psi_{\rm core}/\psi_{\rm new}$ dla punktów z~rodziny rozwiązań.}
\label{tab:psi-ratio}
\begin{tabular}{@{}cccccc@{}}
\toprule
$\alpha^*$ & $a_\Gamma$ & $\lambda^*$ & $\psi_{\rm core}$
    & $\psi_{\rm new} = 1/\sqrt{\lambda^*}$ & $\psi_{\rm core}/\psi_{\rm new}$ \\
\midrule
$5.9148$ & $0.025$ & $2.883\times10^{-6}$ & $1159$ & $589$ & $1.967$ \\
$6.8675$ & $0.030$ & $3.743\times10^{-6}$ & $1009$ & $517$ & $1.952$ \\
$7.7449$ & $0.035$ & $4.618\times10^{-6}$ & $\phantom{0}902$ & $465$ & $1.938$ \\
$8.5616$ & $0.040$ & $5.501\times10^{-6}$ & $\phantom{0}821$ & $426$ & $1.926$ \\
\bottomrule
\end{tabular}
\end{table}

\begin{remark}[Fizyczna interpretacja]%
\label{rem:rdzen-minimum}
Relacja $\psi_{\rm core} \approx 2\,\psi_{\rm new}$ oznacza, że
rdzeń trzeciej generacji leży \emph{za} nowym minimum potencjału
(po przeciwnej stronie bariery~$\psi_{\rm cross}$).
Geometrycznie: rdzeń jest mniej więcej w~połowie drogi między
minimum a~punktem crossover w~skali logarytmicznej.

Analityczne uzasadnienie: warunek samospójności
$g(K_3) = 0$ wymaga równowagi między
\emph{ujemnym} wkładem powłoki ($\psi < \psi_{\rm cross}$)
a~\emph{dodatnim} wkładem rdzenia ($\psi > \psi_{\rm cross}$).
Równowaga ta zachodzi gdy $\psi_{\rm core} \approx \sqrt{3/2}\,\psi_{\rm cross}
= (3/2)\,\psi_{\rm new}$, co jest zgodne z~obserwacją.
\end{remark}

% -------------------------------------------------------------------
\subsection{Wzór Koide'a i~struktura trzech generacji}%
\label{app:I-koide}
% -------------------------------------------------------------------

\begin{proposition}[Wzór Koide'a z~TGP]%
\label{prop:koide}
Dla stosunków mas $1 : r_{21} : r_{31} = 1 : 207 : 3477$
niezmiennik Koide'a~\cite{Koide1983} przyjmuje wartość:
\begin{equation}\label{eq:koide}
    \mathcal{K}
    = \frac{m_1 + m_2 + m_3}{(\sqrt{m_1} + \sqrt{m_2} + \sqrt{m_3})^2}
    = \frac{1 + 207 + 3477}{(1 + \sqrt{207} + \sqrt{3477})^2}
    = 0.66655\ldots,
\end{equation}
co odbiega od wartości $2/3 = 0.66\overline{6}$ o~zaledwie
\begin{equation}
    \frac{\mathcal{K} - 2/3}{2/3} = -0.017\%.
\end{equation}
Wartość eksperymentalna (masy leptonowe PDG~2024):
$\mathcal{K}_{\rm exp} = 0.66666\ldots$ (zgodność $0.006\%$).

Dokładna wartość $\mathcal{K} = 2/3$ przy $r_{21} = 207$
wymaga $r_{31} = 3481$ (różnica od wartości TGP: $+0.11\%$).
\end{proposition}

\begin{remark}[Znaczenie wzoru Koide'a dla struktury generacji]%
\label{rem:koide-znaczenie}
Wzór Koide'a jest relacją między \emph{trzema} argumentami.
Nie istnieje naturalna generalizacja na cztery generacje z~analogicznymi
własnościami algebraicznymi. Jeżeli wzór Koide'a wynika z~głębszej
zasady TGP (np.\ symetria permutacyjna $S_3$ trzech zer $g(K)$),
to \emph{predykuje} dokładnie trzy generacje jako strukturalną
konsekwencję, nie przypadek numeryczny.

Równocześnie: wzór Koide'a dla stosunków $(1, 207, 3477)$
spełniony jest z~dokładnością $0.017\%$ \emph{bez żadnego dopasowania} ---
stosunki te wynikają z~warunku samospójności solitonów TGP.
\end{remark}

% -------------------------------------------------------------------
\subsection{Dlaczego hierarchia zatrzymuje się na trzech generacjach}%
\label{app:I-trzy}
% -------------------------------------------------------------------

Pytanie: dlaczego sekwencja $\{K_1, K_2, K_3\}$ nie kontynuuje się
do $K_4$, $K_5$, \ldots?

\begin{proposition}[Brak czwartej generacji — przyczyna topologiczna]%
\label{prop:brak-czwartej}
Funkcja $g(K)$ dla potencjału~\eqref{eq:V-mod} posiada \emph{dokładnie}
trzy zera. Dla $K > K_3$:
\begin{enumerate}[nosep,label=(\roman*)]
    \item Rdzeń solitonu ma $\psi_{\rm core} > \psi_{\rm cross}$
          (powyżej crossover), zatem wkład rdzenia do $E_{\rm pot}$ jest dodatni.
    \item Powłoka zewnętrzna ($r_{\rm cross} < r < \infty$, gdzie $\psi < \psi_{\rm cross}$)
          daje stały, dominujący \emph{ujemny} wkład do $E_{\rm pot}$,
          rosnący z~$K$ jak $\sim K^2$.
    \item Suma $E(K)/K \to -\infty$ i~$g(K) < 0$ dla wszystkich $K > K_3$ —
          brak czwartego zera.
\end{enumerate}
Potencjał stopnia 8 ($+ \mu\psi^8/8$) dodałby \emph{drugie} przekroczenie
$\psi_{\rm cross,2}$ i~umożliwiłby czwarte zero. Obserwacja trzech
generacji wyklucza taki człon w~potencjale TGP.
\end{proposition}

Dla \textbf{mechanizmu hierarchii tła} $\Phi_0^{(n)}$
(sekwencyjna perturbacja pola przez każdą generację)
zatrzymanie na $n=3$ ma niezależne uzasadnienie:

\begin{remark}[Zatrzymanie hierarchii tła]%
\label{rem:hierarchia-ta}
Sekwencja $\Phi_0^{(1)} \to \Phi_0^{(2)} \to \Phi_0^{(3)}$ jest
\emph{niegeometryczna}: współczynniki wzrostu
$\Phi_0^{(2)}/\Phi_0^{(1)} \approx 2\pi$ i
$\Phi_0^{(3)}/\Phi_0^{(2)} \approx 4.5$ różnią się znacznie
(tabela z~\texttt{badanie\_pi.py}).
Dla bardzo dużego $\Phi_0^{(n)}$ masa ekranowania
$m_{\rm sp} = \sqrt{\gamma/\Phi_0} \to 0$, a~soliton staje się
nieskończenie rozległy — warunek Yukawa przestaje obowiązywać
i~brak zlokalizowanego samospójnego rozwiązania.
W~tym sensie mechanizm tła i~mechanizm kształtu $g(K)$
dają spójną predykcję: \emph{trzy i~tylko trzy generacje}.
\end{remark}

% -------------------------------------------------------------------
\subsection{Wpływ~$\lambda$ na generacje 1~i~2}%
\label{app:I-wplyw-M1-M2}
% -------------------------------------------------------------------

\begin{proposition}[Pomijalny wpływ $\lambda$ na $r_{21}$]%
\label{prop:wplyw-lambda}
Dla pierwszej i~drugiej generacji ($K_1 \ll 1$, $K_2 \approx 2$)
rdzeń solitonu spełnia $\psi_{\rm core} \ll \psi_{\rm new}$
— soliton nie ``widzi'' nowego minimum potencjału.
Korekta członu $\lambda$-owego do $r_{21}$:
\begin{equation}
    \frac{\delta r_{21}}{r_{21}} < 0.03\%,
\end{equation}
co jest poniżej precyzji wyznaczenia $r_{21}$ z~obserwacji.
Mechanizm trzeciej generacji nie zaburza wyznaczonych
wcześniej stosunków mas $M_1:M_2$.
\end{proposition}

% -------------------------------------------------------------------
\subsection{Wyniki numeryczne: status problemów P1--P5}%
\label{app:I-wyniki-P}
% -------------------------------------------------------------------

Poniżej przedstawiamy wyniki analizy numerycznej pięciu problemów
otwartych (skrypty \texttt{p1\_ode\_shooting.py} --
\texttt{p5\_stabilnosc.py}).

% - - - - - - - - - - - - - - - - - - - - - - - - - - - - - -
\subsubsection*{P1 — Weryfikacja ansatzu Yukawa}

Skrypt \texttt{p1\_ode\_shooting.py} całkuje pełne, nieliniowe ODE
dla $\Phi(r)$ (metoda DOP853, rtol $=10^{-10}$), startując od
warunków początkowyh zadanych ansatzem Yukawa w~$r = a_\Gamma = 0.040$.
Wyniki dla kandydatów $K_1$, $K_2$, $K_3$:

\begin{table}[h]
\centering\small
\caption{Wyniki integracji ODE od $r = a_\Gamma$ do $r = 60$.
  Cel: $\Phi(60) = \PhiZero = 1$, $g_{\rm ODE} = 0$.}
\label{tab:p1-wyniki}
\begin{tabular}{@{}ccccc@{}}
\toprule
Kandydat & $K$ & $\Phi_{\rm ODE}(60)$ & $g_{\rm ODE}(K)$ & Wniosek \\
\midrule
$K_1$ & $0.00982$ & $1.00016$ & $-0.314$ & marginalnie OK \\
$K_2$ & $2.03273$ & $0.00000$ & $-3.6\times10^4$ & rozbieżny ($\to 0$) \\
$K_3$ & $34.1445$ & $3036.7\phantom{00}$ & $+2.7\times10^7$ & rozbieżny ($\to +\infty$) \\
\bottomrule
\end{tabular}
\end{table}

\begin{result}[Ansatz Yukawa nie jest IC dla pełnego ODE przy K2, K3]%
\label{res:p1}
Warunki początkowe Yukawa przy $r = a_\Gamma$ prowadzą do rozbieżności
rozwiązania ODE dla $K_2$ i~$K_3$. Oznacza to, że profil Yukawa
\emph{nie} jest dokładnym profilem solitonu dla dużych amplitud
($\psi_{\rm core} \gg 1$).

Prawidłowe shooting powinno startować z~$r = 0$:
$\Phi(0) = \Phi_{\rm core}$ (duże), $\Phi'(0) = 0$
(symetria sferyczna) z~przybliżeniem Taylora
\[
    \Phi''(0) = \frac{V'_{\rm mod}(\Phi_{\rm core})}{3(1 + \alpha/\Phi_{\rm core})}.
\]
Szukanie $\Phi_{\rm core}^*$ takiego że $\Phi(r \to \infty) = \PhiZero$
jest odrębnym zadaniem numerycznym.
Warunek samospójności $g(K_3) = 0$ obliczony w~energii
funkcjonalnej (nie z~ODE) pozostaje główną metodą TGP.
\end{result}

% - - - - - - - - - - - - - - - - - - - - - - - - - - - - - -
\subsubsection*{P2 — Wyznaczenie $\lambda^*$ z~pierwszych zasad}

Skrypt \texttt{p2\_lambda\_mechanizm.py} testuje sześć hipotez (H1--H6)
dla rodziny czterech punktów rozwiązań. Kluczowe wyniki:

\begin{table}[h]
\centering\small
\caption{Test hipotez dla $\lambda^*$.
  Błąd = odchylenie od numerycznego $\lambda^*$.}
\label{tab:p2-hipotezy}
\begin{tabular}{@{}clcc@{}}
\toprule
Hipoteza & Warunek & Typowy błąd & Ocena \\
\midrule
H1 & $|E_{\rm pot,rdzeń}| = |E_{\rm pot,powłoka}|$
   & $\sim 0.85\%$ & nieuniwersalne \\
H2 & $\psi_{\rm core}/\psi_{\rm cross} = \mathrm{const}$
   & wariacja $1.5\%$ & niezachowane \\
H3 & $E_k / |E_p| = 1$ (Virial)
   & $0.6\%$ od jedności & bliskie, nie zero \\
H4 & $r_{31} = 3481$ (Koide dokładne)
   & $\Delta\lambda^*/\lambda^* = 0.23\%$ & bliskie \\
H5 & $\psi_{\rm core} = (e/\!\sqrt{2})\,\psi_{\rm new}$
   & $0.4\text{--}3.0\%$ & najlepsze, niecałkowite \\
H6 & $V_{\rm mod}(\psi_{\rm core})/V_{\rm mod}(\psi_{\rm cross})=\mathrm{const}$
   & wariacja $60\%$ & odrzucone \\
\bottomrule
\end{tabular}
\end{table}

Obserwacja pomocnicza: stosunek $\psi_{\rm cross}/\psi_{\rm new} = 1/\sqrt{3/2}
= 0.8165$ jest \emph{dokładnie} zachowany dla całej rodziny
(wynika algebraicznie z~definicji $\psi_{\rm new}$ i~$\psi_{\rm cross}$).

\begin{result}[$\lambda^*$ jest wolnym parametrem — brak mechanizmu]%
\label{res:p2}
Żadna z~sześciu hipotez nie wyznacza $\lambda^*$ z~błędem $<0.1\%$.
Najlepsza kandydatka (H5: $\psi_{\rm core} = (e/\!\sqrt{2})\,\psi_{\rm new}$)
daje błąd $\lesssim 3\%$, który \emph{maleje} wzdłuż krzywej rozwiązań
do $0.4\%$ w~punkcie $(a_\Gamma = 0.040)$.
Na obecnym etapie $\lambda^*$ pozostaje parametrem wyznaczonym
z~warunku $r_{31} = 3477$ — brak mechanizmu TGP niezależnie
go predykującego.
\end{result}

% ----- Aktualizacja v47 (2026-04-09) -----
\begin{proposition}[$\lambda^* \approx 2/\PhiZero^4$:
  zamkni\k{e}cie numeryczne]%
\label{prop:I-lambda-Z2}
Niech $\PhiZero$ b\k{e}dzie warto\'sci\k{a} t\l{}a kosmologicznego
($\PhiZero = 36\,\Omega_\Lambda \approx 24{,}66$). Relacja
\begin{equation}\label{eq:I-lambda-from-Phi0}
  \boxed{\lambda_{\mathrm{eff}} = \frac{2}{\PhiZero^4}}
\end{equation}
odtwarza $\lambda^*$ ze skanu solitonowego z~odchyleniem
$|\Delta\lambda/\lambda^*| = 1{,}7\%$:
\[
  \frac{2}{\PhiZero^4}\bigg|_{\PhiZero = 24{,}66}
  = 5{,}41 \times 10^{-6}
  \quad\text{vs.}\quad
  \lambda^* = 5{,}50 \times 10^{-6}.
\]

\medskip\noindent\textbf{Motywacja czynnika 2:}
Potencja\l{} $V_{\mathrm{mod}}(\psi) = \psi^3/3 - \psi^4/4
+ \lambda(\psi - 1)^6/6$ stabilizuje pr\'o\.zni\k{e}
z~dwoma sektorami (chiralnym $\mathbb{Z}_2$: kink/antykink).
Ka\.zdy sektor wnosi $1/\PhiZero^4$ do efektywnego
sprz\k{e}\.zenia szostorz\k{e}dowego, daj\k{a}c
\l{}\k{a}czny prefaktor~2.
R\'ownowa\.znie: z~hipotezy $a_\Gamma\cdot\PhiZero = 1$
otrzymujemy $\lambda = 2\,a_\Gamma^4 = 2/\PhiZero^4$.

\medskip\noindent\textbf{Weryfikacja:}
Skrypt \texttt{lambda\_eff\_erg\_vs\_soliton.py}
por\'ownuje sze\'s\'c niezale\.znych \'scie\.zek.
$\lambda = 2/\PhiZero^4$ jest jedyn\k{a} \'scie\.zk\k{a}
z~odchyleniem $< 5\%$.

\medskip\noindent
\textbf{Status:} Propozycja [AN+NUM].
Formalna pochodna czynnika~2 z~$\mathbb{Z}_2$
wymaga potwierdzonej struktury
dw\'och sektor\'ow chiralnych w~ERG ---
program otwarty (OP-$\lambda$).
\end{proposition}

% - - - - - - - - - - - - - - - - - - - - - - - - - - - - - -
\subsubsection*{P4 — Czy $\Phi_0^{(2)} \to 2\pi$ gdy $a_\Gamma \to 0$?}

Skrypt \texttt{p4\_phi0\_limit.py} wyznacza $\Phi_0^{(2)}$ (wartość tła
drugiej generacji) dla $\alpha \in [5.5, 7.0]$ i~$a_\Gamma \in [0.01, 0.10]$,
a~następnie ekstrapoluje do $a_\Gamma \to 0$.

\begin{table}[h]
\centering\small
\caption{Limit $\Phi_0^{(2)}(a_\Gamma \to 0)$ dla różnych $\alpha$.}
\label{tab:p4-phi0}
\begin{tabular}{@{}cccc@{}}
\toprule
$\alpha$ & $\Phi_0^{(2)}(a_\Gamma\to 0)$ & $\Phi_0^{(2)} / 2\pi$ & Odchylenie od $2\pi$ \\
\midrule
$5.5$ & $7.310$ & $1.1635$ & $+16.35\%$ \\
$5.9$ & $7.190$ & $1.1444$ & $+14.44\%$ \\
$6.0$ & $7.163$ & $1.1400$ & $+14.00\%$ \\
$6.5$ & $7.037$ & $1.1199$ & $+11.99\%$ \\
$7.0$ & $6.928$ & $1.1027$ & $+10.27\%$ \\
\bottomrule
\end{tabular}
\end{table}

Dla całego badanego zakresu $\Phi_0^{(2)}(a_\Gamma \to 0) > 2\pi$;
brak zmiany znaku w~$\alpha \in [5, 7]$.
Stosunek $\Phi_0^{(3)}/\Phi_0^{(2)} \approx 3.42\text{--}3.52$
w~limicie $a_\Gamma \to 0$ — bliski $\pi \approx 3.14159$,
ale nie identyczny z~żadną prostą kombinacją stałych matematycznych.

\begin{result}[Hipoteza $\Phi_0^{(2)} \to 2\pi$ obalona w~zakresie $\alpha \in \{5,\ldots,7\}$]%
\label{res:p4}
Ekstrapolacja numeryczna wyklucza $\Phi_0^{(2)} = 2\pi$ dla
$\alpha \leq 7.0$ w~granicy $a_\Gamma \to 0$.
Odchylenie jest zbyt duże ($>10\%$), by wynikało z~artefaktów numerycznych.
Jeśli relacja $\Phi_0^{(2)} = 2\pi$ miałaby istnieć,
wymagałoby to $\alpha > 10$ lub zupełnie innego mechanizmu tła.
Hipoteza ta nie jest poparta przez dostępne obliczenia.
\end{result}

% - - - - - - - - - - - - - - - - - - - - - - - - - - - - - -
\subsubsection*{P5 — Stabilność dynamiczna}

Skrypt \texttt{p5\_stabilnosc.py} dyskretyzuje operator perturbacji
$H = -\nabla^2_r + V''_{\rm mod}(\Phi_{\rm sol}(r))$
na profilu Yukawa i~wyznacza spektrum.

Profil $V''_{\rm mod}(\Phi_{\rm sol}(r))$ wzdłuż solitonu $K_3$:
\begin{itemize}[nosep]
\item $r = a_\Gamma$: $V'' = +1.04 \times 10^7$ (rdzeń, silnie stabilizujący)
\item $r \approx 0.091$: $V'' = 0$ (zmiana znaku)
\item $r > 0.1$: $V'' < 0$, dążąc do $V''(1) = -1$ przy $r \to \infty$
\end{itemize}

Najniższe wartości własne $H$ dla $K_3$:
$E_0 \approx -8.0 \times 10^4$, $E_1 \approx -7.1\times10^4$, \ldots
(wszystkie $\ll V''(1) = -1$).

\begin{result}[Wstępna analiza wskazuje niestabilność — wyniki niepełne]%
\label{res:p5}
Preliminarna dyskretyzacja daje wszystkie mody niestabilne
($E_i \ll -1$). Interpretacja wymaga ostrożności z~trzech powodów:
\begin{enumerate}[nosep,label=(\roman*)]
\item profil solitonu jest ansatzem Yukawa, nie rozwiązaniem dokładnym ODE;
\item operator kinetyczny nie uwzględnia członu $(1 + \alpha/\Phi)$ z~lagranżianu TGP;
\item tło $\Phi = \PhiZero$ jest \emph{siodłem} potencjału ($V''(1) = -1 < 0$):
      continuum perturbacji tła samo w~sobie jest niestabilne.
\end{enumerate}
Test Derricka potwierdza brak skalowalnej równowagi:
$-E_k + 3E_p \neq 0$ dla wszystkich trzech solitonów.
Pełna analiza stabilności wymaga rozwiązania P1 (dokładny profil ODE).
\end{result}

% -------------------------------------------------------------------
\subsection{Problemy otwarte (zaktualizowane)}%
\label{app:I-otwarte}
% -------------------------------------------------------------------

\begin{problem}[Shooting od $r=0$: właściwe równanie ODE dla solitonów]%
\label{prob:ODE}
Wynik P1 (tabela~\ref{tab:p1-wyniki}) pokazuje, że ansatz Yukawa
nie jest właściwym warunkiem początkowym dla full-ODE.
Konieczna jest implementacja shootingu od $r = 0$:
\begin{equation}
    \Phi''(0) = \frac{V'_{\rm mod}(\Phi_0^{\rm try})}{3(1 + \alpha/\Phi_0^{\rm try})},
    \quad \Phi'(0) = 0,
\end{equation}
i~wyznaczenie $\Phi_0^{\rm try}$ takie, że $\Phi(r) \to 1$ dla $r \to \infty$.
\textbf{Pytanie otwarte}: czy solitony $M_2$ i~$M_3$ istnieją jako
dokładne statyczne rozwiązania ODE, czy tylko jako stacjonarności funkcjonału energii?
\end{problem}

\begin{problem}[Wyznaczenie $\lambda^*$ z~pierwszych zasad]%
\label{prob:lambda-star}
Wynik P2 (tabela~\ref{tab:p2-hipotezy}) wskazuje, że $\lambda^*$
nie jest wyznaczony przez żadną z~testowanych relacji.
Najbardziej obiecująca kandydatka H5 ($\psi_{\rm core} = (e/\!\sqrt{2})\,\psi_{\rm new}$)
ma błąd $0.4\text{--}3\%$ — zbyt duży na mechanizm zasadniczy.

\textbf{Pytanie otwarte}: czy $\lambda$ wynika z~głębszej symetrii TGP,
czy jest nowym niezależnym parametrem teorii (analogicznym do stałej
sprzężenia w~QED)?
\end{problem}

\begin{problem}[Jednoparametrowa rodzina — fizyczny selektor]%
\label{prob:rodzina}
Rodzina rozwiązań $(\alpha^*, a_\Gamma, \lambda^*)$ tworzy krzywą
(tabela~\ref{tab:rodzina-rozwiazan}).
Każdy punkt daje $r_{21} = 207$, $r_{31} = 3477$.

\textbf{Obserwacja (wynik P2)}: stosunek $\psi_{\rm core}/\psi_{\rm cross}$
maleje monotonicz\-nie wzdłuż krzywej:

\medskip
\begin{center}\small
\begin{tabular}{@{}cccc@{}}
\toprule
$a_\Gamma$ & $\psi_{\rm core}/\psi_{\rm cross}$ & $\pi/2 = 1.5708$ & Odchylenie \\
\midrule
$0.025$ & $1.6062$ & & $2.25\%$ \\
$0.030$ & $1.5937$ & & $1.46\%$ \\
$0.035$ & $1.5826$ & & $0.75\%$ \\
$0.040$ & $1.5726$ & $\leftarrow$ & $0.11\%$ \\
\bottomrule
\end{tabular}
\end{center}
\medskip

Punkt $a_\Gamma = 0.040$ jest najbliżej relacji $\psi_{\rm core}/\psi_{\rm cross} = \pi/2$,
co odpowiada
\begin{equation}\label{eq:selektor-pi}
    \psi_{\rm core} = \frac{\pi}{2}\,\psi_{\rm cross}
    = \frac{\pi}{2}\sqrt{\frac{3}{2}}\,\psi_{\rm new}
    = \pi\sqrt{\frac{3}{8}}\;\psi_{\rm new}
    \approx 1.924\,\psi_{\rm new}.
\end{equation}
Obserwowana wartość dla $a_\Gamma=0.040$: $\psi_{\rm core}/\psi_{\rm new} = 1.926$
(odchylenie od~\eqref{eq:selektor-pi}: $+0.11\%$).
Ekstrapolacja liniowa daje $a_\Gamma^* \approx 0.041$ dla dokładnej równości.

\textbf{Hipoteza P3}: fizycznym selektorem rodziny jest warunek
$\psi_{\rm core} = (\pi/2)\,\psi_{\rm cross}$,
czyli ``rdzeń solitonu leży w~ćwiartce drogi ponad punktem crossover''
w~skali mnożnikowej.

\textbf{Pytanie otwarte}: skąd wynika czynnik $\pi/2$?
Możliwe uzasadnienia:
(a) geometria całki sferycznej $\int_0^{\pi/2} \sin^2\theta\,d\theta = \pi/4$;
(b) warunek stacjonarności fazy profilu Yukawa $e^{-r_{\rm cross}}/r_{\rm cross}$;
(c) przypadek numeryczny wymagający weryfikacji w~szerszej rodzinie parametrów.
\end{problem}

\begin{problem}[Stabilność dynamiczna solitonów]%
\label{prob:stabilnosc-M3}
Wynik P5 jest wstępny i~niedokładny (patrz uwagi w~\ref{res:p5}).
Kluczowym krokiem jest uzyskanie dokładnego profilu ODE z~P1,
a~następnie analiza spektrum z~pełnym operatorem kinetycznym.

\textbf{Ważna obserwacja}: $V''_{\rm mod}(1) = -1 < 0$ oznacza, że
tło $\Phi = \PhiZero$ jest siodłem — nie minimum globalnym.
Stabilność solitonu musi być rozumiana \emph{względem tła},
nie absolutnie. Długoczasowa dynamika wymaga uwzględnienia
rozpadu tła na $\psi_{\rm new} \approx 426$ lub mechani\-zmu
stabilizacji (np.\ ładunek topologiczny, bariera potencjału).
\end{problem}

% -------------------------------------------------------------------
\subsection{Podsumowanie wyników}%
\label{app:I-podsumowanie}
% -------------------------------------------------------------------

\begin{table}[h]
\centering
\small
\caption{Zestawienie wyników modyfikacji potencjału TGP
  $V_{\rm mod} = V_{\rm orig} + \lambda(\psi-1)^6/6$.}
\label{tab:podsumowanie-v2}
\begin{tabular}{@{}p{6.5cm}c p{4.5cm}@{}}
\toprule
\textbf{Twierdzenie / wynik} & \textbf{Status} & \textbf{Skrypt} \\
\midrule
Warunki próżni zachowane & Dowiedziony analitycznie
    & Stwierdzenie~\ref{prop:proznia-zachowana} \\
Trzy zera $g(K)$ dla stopnia 6 & Dowiedziony analitycznie
    & Tw.~\ref{thm:trzy-zera} \\
Poprzednia ``rozbieżność'' = artefakt siatki
    & Potwierdzony num.
    & \texttt{v2\_kwadryatura.py} \\
Trzecie zero rzeczywiste (zbieżne)
    & Potwierdzony num.
    & \texttt{v2\_konwergencja\_finalna.py} \\
$r_{21} = 207$, $r_{31} = 3477$ jednocześnie
    & Potwierdzony num.
    & \texttt{v2\_bisekcja\_alpha.py} \\
Rodzina jednoparametrowa (4 punkty)
    & Potwierdzony num.
    & Tabela~\ref{tab:rodzina-rozwiazan} \\
$\psi_{\rm core}/\psi_{\rm new} \approx 2$ (universalne)
    & Obserwacja num.
    & Tabela~\ref{tab:psi-ratio} \\
Wzór Koide'a spełniony do $0.017\%$
    & Weryfikacja num.
    & \texttt{badanie\_pi.py} \\
Brak czwartej generacji (predykcja)
    & Analitycznie + num.
    & Stwierdzenie~\ref{prop:brak-czwartej} \\
\midrule
Ansatz Yukawa = IC dla full-ODE? (P1)
    & \textbf{Nie} dla $K_2$, $K_3$
    & Wynik~\ref{res:p1}, tab.~\ref{tab:p1-wyniki} \\
$\lambda^*$ z~pierwszych zasad (P2)
    & \textbf{Brak mechanizmu}
    & Wynik~\ref{res:p2}, tab.~\ref{tab:p2-hipotezy} \\
$\Phi_0^{(2)} \to 2\pi$ dla $\alpha \leq 7$ (P4)
    & \textbf{Obalona} ($>10\%$ odchyleń)
    & Wynik~\ref{res:p4}, tab.~\ref{tab:p4-phi0} \\
Stabilność dynamiczna $M_3$ (P5)
    & \textbf{Wstępna niestabilność} (niepełna analiza)
    & Wynik~\ref{res:p5} \\
Selektor punktu na krzywej rozwiązań (P3)
    & \textbf{Otwarty}
    & Problem~\ref{prob:rodzina} \\
Shooting ODE od $r=0$ (P1 rozszerzony)
    & \textbf{Otwarty}
    & Problem~\ref{prob:ODE} \\
\bottomrule
\end{tabular}
\end{table}
